% =============================================================
%  Chapter 7 — Conclusion
% =============================================================

\chapter{Conclusion}
\label{ch:conclusion}


% ----- §7.1 Thesis Summary -----
\section{Thesis Summary}
\label{sec:conclusion-summary}

\textsc{BaxBench}~\cite{vero2025baxbench} and similar benchmarks established
that LLMs can generate functionally correct backend applications, but they
evaluate every candidate in an isolated, single-machine Docker container
and stop at functional and security correctness
(Section~\ref{sec:intro-context}). Nothing in that evaluation methodology
says anything about whether the generated application would serve
production traffic well once it is actually deployed: replica counts,
resource requests, pod placement, and database topology are all decided
outside the scope of the benchmark, by a human operator or a reactive
autoscaler that reasons about none of it jointly
(Section~\ref{sec:intro-gap}). This thesis asked whether an LLM agent
could close that gap itself: iteratively optimise the deployment
configuration of a backend application on a Kubernetes cluster, guided by
runtime performance feedback, to maximise sustainable throughput under
fixed infrastructure constraints (Section~\ref{sec:intro-missing}).

The approach is an iterative benchmarking framework that extends
\textsc{BaxBench} with a real multi-node Kubernetes cluster and a
five-stage refinement loop --- decision, code, spec, deploy, bench
(Chapter~\ref{ch:methods}) --- in which an LLM agent proposes exactly one
change per iteration, to either application code or a constrained
deployment specification (\texttt{spec.yaml}), and receives structured
load-test and resource-utilisation feedback before the next iteration.
Performance is scored by \emph{sustained goodput} under an adaptive load
profile that discovers each deployment's sustainable throughput without
per-scenario rate tuning (Section~\ref{sec:methods-load-profile}). The
methodological contributions this required --- a schema that bounds the
agent's deployment search space to meaningful decisions rather than raw
Kubernetes YAML, a failure taxonomy and routing scheme that keeps a
single bad iteration from silently corrupting the next one
(Section~\ref{sec:methods-failure}), and an evaluation harness that
tracks LLM cost as a first-class quantity alongside goodput
(Section~\ref{sec:methods-workflow}) --- are what make the empirical
results in Chapter~\ref{ch:results} attributable to specific causes
rather than an unexplained end-to-end number.


% ----- §7.2 Key Findings -----
\section{Key Findings}
\label{sec:conclusion-findings}

\paragraph{What worked.} Iterative refinement improved sustained goodput
in essentially every cell that ever became benchmarkable, across all
three evaluated models, and in a substantial fraction of cells it did not
just improve an already-working deployment but converted a non-functional
one into a functional one (Section~\ref{sec:results-rq1}). Where a
trajectory ran long enough, the agent's deployment-specification choices
tracked the workload's actual scaling characteristics --- favouring
horizontal fan-out over per-pod headroom for a near-stateless workload,
without ever being told that heuristic (Section~\ref{sec:results-rq5}).
Keeping code and deployment refinement mutually exclusive per iteration
turned out to be more than a bookkeeping convenience: the two levers
measurably specialise in practice, with code refinement carrying almost
all of the largest single-step recoveries and deployment refinement
acting as a smaller, more consistent capacity adjustment
(Section~\ref{sec:results-rq2}).

\paragraph{What did not work.} Deployment refinement, taken on its own,
was occasionally counterproductive: a well-reasoned response to correct
diagnostics regressed a healthy deployment to zero goodput, because
nothing prevents one refinement step from changing several
deployment-spec fields whose interaction it cannot fully anticipate
(Section~\ref{sec:results-rq6}). Separately, and closer to a framework
finding than a model one: static spec validation did not catch every
deployment that would later prove impossible to schedule, producing this
evaluation's first (and only) deploy-phase failures
(Section~\ref{sec:results-rq6}).

\paragraph{What surprised us.} The multiplicative gain from refinement
turned out to depend much more on how badly under-provisioned the
baseline happened to be than on any consistent ``value added by
refinement'' figure: the same framework produced gains ranging from
just over $1\times$ to more than $75\times$ depending on the cell
(Section~\ref{sec:results-rq1}). The largest surprise, however, was
methodological rather than empirical: what first looked like a clean
model-capability signal (\texttt{gpt-5.5} completing every cell,
\texttt{claude-opus-4-8} completing noticeably fewer) turned out, once
the underlying error text was actually read cell by cell, to be entirely
explained by one provider account running low on funds during data
collection, not by generation quality (Section~\ref{sec:results-rq3}).
Every cell that never reached a baseline in this dataset shows the same
handful of billing/quota error strings, not a functional-test failure,
and every one of them belongs to \texttt{claude-opus-4-8}: re-tracing the
same check against \texttt{glm-5.2}, which had looked similarly affected
in an earlier pass over an earlier snapshot of this dataset, found its
completion rate fully recovered (12/12, matching \texttt{gpt-5.5}) once
that snapshot was re-aggregated, with zero billing-kind failures
remaining anywhere in its data. That this was not obvious from the
aggregate failure counts alone, since a coarse \texttt{llm\_call} label
conflated it with genuine transport faults until individual records were
inspected, cell by cell, snapshot by snapshot, is itself a finding
about how easy it is to over-read a completion funnel in agentic LLM
evaluation. Also unexpected, though smaller: the deployment-and-cluster
path was \emph{nearly}, but not quite, invisible as a source of failure.
Four spec-refinement iterations, on two cells from two different
models, produced deployments that passed static validation but could not
be deployed as intended, the only deploy-phase failures recorded anywhere
in this evaluation (Section~\ref{sec:results-rq6}); a further, entirely
new, framework-specific correctness signal (\texttt{glm-5.2}'s repeated
\textsc{Rust-Actix} compile failures) only became visible once its cells
were no longer cut short by the billing confound
(Section~\ref{sec:results-rq4}).

\paragraph{What we learned from the evaluation.} The single biggest
determinant of final performance in this dataset was not deployment
tuning at all, but whether the generated code was correct enough to
survive real load in the first place, including failure modes (a
connection pooler's transaction-mode incompatibility with server-side
prepared statements) that a single-container functional-test suite
structurally cannot observe (Section~\ref{sec:results-rq6}). Evaluating a
deployment-optimisation agent in isolation from code correctness would
have missed this; the two are entangled in practice, even though the
framework deliberately separates them as levers.


% ----- §7.3 Research Question: Final Answer -----
\section{Research Question: Final Answer}
\label{sec:conclusion-rq}

Restating the thesis's central question (Section~\ref{sec:intro-missing}):
\emph{can an LLM agent iteratively optimise the deployment configuration
of a backend application on a Kubernetes cluster, guided by runtime
performance feedback, to maximise sustainable throughput under fixed
infrastructure constraints?}

\textbf{Yes, within the scope this thesis tested.} Across 36 scenario
$\times$ framework cells and three models, an LLM agent that alternates
between code and deployment-specification refinement, guided only by
adaptive load-test feedback and resource diagnostics, consistently
increased sustained goodput once a functionally correct baseline
existed, and in a majority of dead-baseline cells it produced a working
deployment where none existed initially. The deployment-specification
lever specifically --- the part of the question closest to this thesis's
actual contribution, since code-level LLM optimisation is comparatively
well studied (Section~\ref{sec:rw-llm-perf}) --- behaved as a genuine,
if bounded, optimisation signal: spec trajectories moved toward
workload-appropriate configurations without being given target ranges
(Section~\ref{sec:results-rq5}), and static validation plus the
framework's failure routing kept bad deployment decisions from silently
compounding (Section~\ref{sec:results-rq6}).

The qualification matters as much as the answer. What this thesis cannot
yet say is whether the result generalises beyond four database-backed
scenarios, one cluster topology, and a ten-iteration budget at $n=1$
sample per cell (Section~\ref{sec:conclusion-limitations}); nor can it
cleanly separate ``the agent optimises deployment configuration well''
from ``the agent optimises deployment configuration well \emph{once code
correctness stops being the bottleneck}'' --- in this dataset, the two
were rarely simultaneously true early in a trajectory.


% ----- §7.4 Limitations -----
\section{Limitations}
\label{sec:conclusion-limitations}

The full discussion is in Section~\ref{sec:discussion-limitations}; in
summary, every result in this thesis should be read against: six of 36
cells interrupted by provider billing/quota exhaustion during data
collection, four of them never reaching a working baseline as a result,
all six attributable to \texttt{claude-opus-4-8} or \texttt{gpt-5.5},
which makes the model- and framework-comparison findings (RQ3, RQ4)
weaker specifically where \texttt{claude-opus-4-8}'s \textsc{Rust-Actix}
cells are involved, while leaving the \texttt{gpt-5.5}-vs-\texttt{glm-5.2}
comparison and the code/deployment-refinement findings (RQ1, RQ2, RQ6)
largely unaffected; scenario coverage limited to 4 of \textsc{BaxBench}'s 28
scenarios, all database-backed; a fixed 10-iteration refinement budget;
$n=1$ sample per cell, so no comparative claim here has been tested
against sampling variance; a single scalar (sustained goodput under one
adaptive load profile) as the sole optimisation signal; temperature-$0.2$
but still stochastic LLM behaviour with no repeated draws to quantify it;
and a single cluster topology and hardware profile, so ``fixed
infrastructure constraints'' in this thesis means one specific capacity
envelope, not infrastructure constraints in general.


% ----- §7.5 Future Work -----
\section{Future Work}
\label{sec:conclusion-future-work}

\paragraph{Re-run the remaining billing-interrupted cells.} A previous
pass over an earlier snapshot of this dataset identified thirteen cells
across all three models needing a re-run; re-aggregating against a newer
snapshot for this thesis shows \texttt{glm-5.2}'s six are now fully
resolved (12/12 completion, zero billing-kind failures), narrowing the
open item to \texttt{claude-opus-4-8}. The single highest-value next step
is mechanical, not conceptual: re-run \texttt{claude-opus-4-8}'s four
non-baseline \textsc{Rust-Actix} cells, and ideally its one truncated
\textsc{Recipes}$\times$\textsc{Python-Flask} cell and \texttt{gpt-5.5}'s
one truncated \textsc{Recipes}$\times$\textsc{Rust-Actix} cell, with a
provider account funded for the full trajectory, so that RQ3's model
comparison and RQ4's \textsc{Rust-Actix} framework comparison can include
\texttt{claude-opus-4-8} rather than leaving it open. Beyond recovering
those specific cells, this argues for a concrete process change for
future runs of this framework: track remaining provider balance as a
first-class precondition the harness checks before starting a cell
(analogous to the existing \texttt{--llm-max-cost} per-experiment cap,
Section~\ref{sec:methods-workflow}), rather than discovering the account
ran out partway through a ten-iteration trajectory. Once the account-level
gap is closed, add repeated samples per cell so that the ratios in
Chapter~\ref{ch:results} can be reported with variance rather than as
single draws.

\paragraph{Broaden scenario coverage beyond database-bound workloads.}
Every scenario evaluated here is network- and database-bound
(Section~\ref{sec:discussion-limitations}); extending to compute-bound or
memory-bound \textsc{BaxBench} scenarios would test whether the
code/deployment division of labour observed here (Section~\ref{sec:results-rq2})
is a general property of the two levers or an artefact of the scenarios
chosen.

\paragraph{Constrain or measure within-lever attribution.} The backfire
case in Section~\ref{sec:results-rq6} shows that a single spec-refinement
step can change several fields at once, which limits attribution the
same way mixing code and spec changes would. Future work could either
constrain spec-refinement steps to one field at a time and measure
whether that improves reliability, or instrument enough to attribute
outcomes to individual fields \emph{post hoc} without constraining the
agent's search space.

\paragraph{Test transfer across cluster topology and scale.} All results
come from one cluster profile (Section~\ref{sec:setup-infrastructure}).
Repeating even a subset of cells on a cluster an order of magnitude
larger or smaller, or with a different node/network topology, would show
whether the deployment decisions the agent converges to
(Section~\ref{sec:results-rq5}) are tuned to the workload or partly to
this specific cluster's capacity envelope.

\paragraph{Richer feedback and multi-objective scoring.} Sustained
goodput is deliberately a single number
(Section~\ref{sec:methods-load-profile}); resource-utilisation and
per-endpoint latency diagnostics are already collected
(Section~\ref{sec:methods-outcome}) but not scored. Feeding
cost-per-request or tail-latency objectives into the decision stage
alongside goodput --- rather than only into the human-readable feedback
--- is a natural extension once the single-objective case is on firmer
statistical footing.
